% !PRINTED 588
\weberpage{588}

Da $A$ und $D$ ohne gemeinschaftliche Teiler sind, so ist für eine Variable $x$ und eine Konstante $h$:
\[
h\,x^{m-1} A\!\left(\frac{1}{x}\right) = D(x),
\]
und indem man $x$ durch $1/x$ ersetzt:
\[
h\,A(x) = x^{m-1} D\!\left(\frac{1}{x}\right),
\]
und nach (16):
\[
h^2 = (-1)^{cx}.
\]
Also ist
\weq{17}{D(x) = \varepsilon\,x^{m-1} A\!\left(\frac{1}{x}\right),}
worin $\varepsilon = \pm 1$ oder $= \pm i$ ist. Das hängt nach (16) davon ab, ob $cx$ gerade oder ungerade ist\footnote{Aus § 138 ergibt sich, daß $i$ im Körper $\mathfrak{L}$ enthalten ist; (17) steht also nicht im Widerspruch mit dem Satze 2.}.

Gebrauchen wir also wieder die Bezeichnung (10), jedoch jetzt unter der Voraussetzung, daß $A$ und $D$ ohne gemeinschaftliche Teiler sind, so ergibt sich:
\weq{18}{\varepsilon\,\boldsymbol{s}(\mu u) = \frac{A_1\boldsymbol{s} + A_3\boldsymbol{s}^3 + \cdots A_{m-2}\boldsymbol{s}^{m-2} + \boldsymbol{s}^m}{1 + A_{m-2}\boldsymbol{s}^2 + \cdots A_3\boldsymbol{s}^{m-3} + A_1\boldsymbol{s}^{m-1}},}
und die $A_1, A_3, \ldots, A_{m-1}$ sind \gesperrt{ganze Zahlen des Körpers} $\mathfrak{L}$, denn sie sind die symmetrischen Grundfunktionen der ganzen Zahlen (13).
% [source sic: the prose reaches A_{m-1}, whereas equation (18) ends the unknown sequence at A_{m-2}]

Im besonderen ergibt sich, wenn man $u=0$ setzt:
\weq{19}{A_1 = \varepsilon\mu,}
und daraus:
\weq{20}{\varepsilon\mu = \prod\nolimits^{\varrho} \boldsymbol{s}\!\left(\frac{2\varrho}{\mu}\right).}
Die $\boldsymbol{s}\!\left(\dfrac{2\varrho}{\mu}\right)$ sind also Teiler von $\mu$.

Aus (10) folgt:
\weq{21}{\boldsymbol{s}A(\boldsymbol{s}) - \boldsymbol{s}(\mu u)D(\boldsymbol{s}) = 0,}
und dies ist, wenn $\boldsymbol{s}(\mu u)$ als gegeben betrachtet wird, eine Gleichung für $\boldsymbol{s}(u)$ vom $m^{\text{ten}}$ Grade, deren Wurzeln sind:
\[
\boldsymbol{s}(u),\ \boldsymbol{s}\!\left(u+\frac{2\varrho}{\mu}\right).
\]
%CONT>>
